\documentclass{ximera}
%% Macros shared by every activity. Each activity loads this file with
%% %% Macros shared by every activity. Each activity loads this file with
%% %% Macros shared by every activity. Each activity loads this file with
%% \input{../preamble.tex}, and the xourse file with \input{preamble.tex}.

%% MathML intent for screen readers: \xmunits, \xmabs, \xmpoint and the
%% other commands of TailorSwiftBot.sty. The file is latex/xmintent.sty.
\usepackage{xmintent}

%% The interactive pieces of Ximera: \answer, multipleChoice and \choice,
%% selectAll, hint and the rest. Keep this line: without it those commands
%% are undefined and come out as plain text. The file is
%% latex/ximera-components.sty.
\usepackage{ximera-components}

%% Tables with the tabularray package. The settings rowhead, xmcolhead,
%% caption and summary make their header cells known to screen readers, on
%% the web page and in the tagged PDF.
\usepackage{tabularray}

\newcommand{\RR}{\mathbb{R}}
\newcommand{\dfn}{\textit}
, and the xourse file with \input{preamble.tex}.

%% MathML intent for screen readers: \xmunits, \xmabs, \xmpoint and the
%% other commands of TailorSwiftBot.sty. The file is latex/xmintent.sty.
\usepackage{xmintent}

%% The interactive pieces of Ximera: \answer, multipleChoice and \choice,
%% selectAll, hint and the rest. Keep this line: without it those commands
%% are undefined and come out as plain text. The file is
%% latex/ximera-components.sty.
\usepackage{ximera-components}

%% Tables with the tabularray package. The settings rowhead, xmcolhead,
%% caption and summary make their header cells known to screen readers, on
%% the web page and in the tagged PDF.
\usepackage{tabularray}

\newcommand{\RR}{\mathbb{R}}
\newcommand{\dfn}{\textit}
, and the xourse file with \input{preamble.tex}.

%% MathML intent for screen readers: \xmunits, \xmabs, \xmpoint and the
%% other commands of TailorSwiftBot.sty. The file is latex/xmintent.sty.
\usepackage{xmintent}

%% The interactive pieces of Ximera: \answer, multipleChoice and \choice,
%% selectAll, hint and the rest. Keep this line: without it those commands
%% are undefined and come out as plain text. The file is
%% latex/ximera-components.sty.
\usepackage{ximera-components}

%% Tables with the tabularray package. The settings rowhead, xmcolhead,
%% caption and summary make their header cells known to screen readers, on
%% the web page and in the tagged PDF.
\usepackage{tabularray}

\newcommand{\RR}{\mathbb{R}}
\newcommand{\dfn}{\textit}


\xmnewfunction{\F}{\mathcal{F}}
\xmnewfunction[Fourier-transform]{\FT}{\mathcal{F}}

\title{Math Intent}

\begin{document}
\begin{abstract}
Testing MathML intent for screen readers.
\end{abstract}
\maketitle

Each formula below should be read by NVDA with MathCAT as described.

An absolute value, read ``the absolute value of x minus 3'': $\xmabs{x-3}$.

A unit, read ``5 inches'': $5\,\xmunits{in}$.

A unit with a power, read ``130 centimeters cubed'':
\[
V = 130\,\xmunits{cm}^3 .
\]

A compound unit, read ``10 meters divided by seconds'': $10\,\xmunits{m}/\xmunits{s}$.

A unit with a name of its own, read ``2 square feet'': $2\,\xmunits[square-feet]{sq.~ft.}$

For comparison, the same bars with no intent: $|x-3|$.

\section*{More commands}

In each line the reading with intent is given first.

A determinant, ``the determinant of A'': $\xmdet{A}$.

A point, ``the point 2 comma 5'': $\xmpoint{2,5}$. In space: $\xmpoint{1,-2,z}$.

Intervals: closed $\xminterval{[a,b]}$, open $\xminterval{(0,1)}$,
closed open $\xminterval{[0,\infty)}$, open closed $\xminterval{(-1,3]}$.

A function made with \verb|\xmnewfunction|, ``script F of f of t'':
$\F[f](t)$. With a spoken name, ``Fourier transform of f of xi'':
$\FT[f](\xi)$. Alone: $\F$.

A one-off function, ``R of x'': $\xmfunction{R}(x)$; two variables:
$\xmfunction{g}(x,y)$.

Composition, ``f composed with g of x'': $\xmcompose{f}{g}(x)$; alone:
$\xmcompose{f}{g}$.

An inverse, ``f inverse of x'': $\xminv{f}(x)$; a matrix: $\xminv{A}$.

A transpose: $\xmtrans{A}$.

Vectors: $\xmvec{v}$, $\xmvec*{v}$, unit vectors $\xmunitvec{u}$,
$\xmunitvec*{u}$.

Products: $\xmdot{\xmvec{u}}{\xmvec{v}}$ and
$\xmcross{\xmvec{u}}{\xmvec{v}}$.

A norm: $\xmnorm{\xmvec{v}}$; a magnitude: $\xmnorm[magnitude]{\xmvec{v}}$.

A cardinality: $\xmcard{A}$.

Degrees: $90\xmdeg$. Outside math: a right angle is 90\xmdeg, and the
board is 5~\xmunits{ft.} long.

A blank to fill in. Name: \xmblank[2.5in] \quad Date: \xmblank

In display math:
\[
\xmabs{\xmdot{\xmvec{u}}{\xmvec{v}}} \le \xmnorm{\xmvec{u}}\,\xmnorm{\xmvec{v}},
\qquad \xmdet{\xmtrans{A}} = \xmdet{A} .
\]

\begin{problem}
An answer box next to a unit: the square has side $\answer{4}\,\xmunits{cm}$
when its area is $16\,\xmunits{cm}^2$.
\end{problem}

\end{document}
